\section{Responsible Scaling Policy（RSP）}

\subsection{RSP 的核心框架}

\begin{importantbox}{ASL（AI Safety Level）分级}
Anthropic 提出 AI Safety Level 分级系统（类似生物安全等级 BSL）：
\begin{itemize}
    \item \textbf{ASL-1}：无显著风险的 AI 系统
    \item \textbf{ASL-2}：当前大模型水平，需要标准安全措施
    \item \textbf{ASL-3}：具备显著自主能力或专业知识，需要强化安全措施
    \item \textbf{ASL-4+}：接近超人能力，需要极端安全措施
\end{itemize}
每个级别对应不同的安全要求——部署条件、监控强度、使用限制。
\end{importantbox}

\subsection{RSP 的运作机制}

\begin{enumerate}
    \item \textbf{定期评估}：每次训练新模型或显著更新时，评估其能力是否跨越 ASL 边界
    \item \textbf{触发机制}：如果评估显示模型接近下一级能力，暂停部署直到相应安全措施就位
    \item \textbf{安全投入}：将安全研发预算与模型能力增长挂钩
    \item \textbf{透明度}：定期发布能力评估报告
\end{enumerate}

\subsection{RSP 的意义}

\begin{knowledgebox}{为什么 RSP 重要}
RSP 的核心贡献是将``模型能力''与``安全措施''显式绑定——不是``先发布再补救''，而是``先确保安全再推进''。这为整个行业提供了一个可操作的安全框架模板。
\end{knowledgebox}

\subsection{本章小结}

RSP 通过 ASL 分级和触发机制，实现了``能力驱动的安全升级''。这是目前最具操作性的 AI 安全框架之一。

